\documentclass[10pt,letterpaper]{article}
\usepackage[margin=12mm]{geometry}
\usepackage{graphicx}
\usepackage{caption}
\pagestyle{empty}
\captionsetup{font=small,labelfont=bf}
\begin{document}
\begin{figure}[p]
\centering
\includegraphics[width=\textwidth]{output/tracking_s3_amendment/figS3C_G_ce_tracking_variance.pdf}
\caption*{\textbf{Figure S3 (continued). Assignment stability, transfer costs, and common improvements across tracking replicates.}
All panels use the same 275 matched terminal cells, parent capacities, and fixed
top-20 protein representation as panels A--B. Optimization uses 301
endpoint-inclusive weights; travel weight 0 minimizes cell-state cost and
weight 1 minimizes travel. Weighted sweeps scale travel and cell-state costs
separately by exact first-cousin-null standard deviations.
\textbf{(C)} Agreement in biological parent identity at corresponding weights.
Exchanging slots of the same parent is not an edge change. Colors identify
replicate pairs in C--F. Correspondence by canonical front position gives
similar median agreement (76.0--79.6\%).
\textbf{(D)} Mean overlap of each child's nearest distinct candidate-parent
sets. The dotted line is the expectation for independent uniformly chosen sets
from 172 parents; it is not a spatial null or significance test.
\textbf{(E)} Excess target-geometry weighted cost of a transferred source
assignment, divided by the target optimum at the same weight. Solid and dashed
curves distinguish transfer directions. This is an optimization diagnostic,
not a biological combined-efficiency score.
\textbf{(F)} At the pure-travel endpoint, excess target travel divided by the
target's natural-to-optimal travel saving. Values at or above 100\% erase that
saving; this is a travel-axis statement, not two-objective dominance by nature.
\textbf{(G)} Joint integer optimization maximizes the smallest percentage travel
reduction across embryos while requiring at least 0\%, 0.5\%, or 1\% reduction
in shared cell-state cost. Each group of three markers evaluates one common
assignment; the x-coordinate shows its achieved cell-state reduction. Dashed
segments guide the eye and are not additional feasible assignments. All three
solutions improve both objectives in all three geometries and are certified
optimal within numerical solver tolerances under the stated minimax criterion.
They were fitted to these three embryos; held-out generalization is untested.
Replicate differences combine biological variation, tracking error, and residual
stage differences. No population uncertainty or isolated tracking-error variance
is inferred.}
\end{figure}
\end{document}
